\documentclass{article}
\usepackage{amsmath,amssymb,amsthm}
\usepackage{algorithm}
\usepackage{algpseudocode}
\usepackage{hyperref}
\usepackage{enumitem}
\usepackage{xcolor}
\newcommand{\E}{\mathbb{E}}
\newcommand{\R}{\mathbb{R}}
\newcommand{\norm}[1]{\left\lVert#1\right\rVert}
\newtheorem{lemma}{Lemma}
\newtheorem{theorem}{Theorem}
\newtheorem{assumption}{Assumption}

\begin{document}

\section*{Algorithm Name}
\textbf{Stochastic Variance‑Reduced Gradient (SVRG) with Double Loop}\\
The method alternates between a \emph{snapshot} phase (outer loop) where a full gradient is computed, and an \emph{inner} phase that performs $m$ stochastic updates using a variance‑reduced estimator.

\section*{Mathematical Formulation}
Let the finite‑sum objective be
\[
f(x)\;=\;\frac{1}{n}\sum_{i=1}^{n} f_i(x),\qquad x\in\R^{d},
\]
where each $f_i$ is $L$‑smooth.  
Denote by $\tilde{x}^{s}$ the snapshot point at outer iteration $s$ and by $x_{t}^{s}$ the iterate after $t$ inner steps of epoch $s$.
The algorithm proceeds as follows:

\begin{enumerate}[label=\arabic*.]
\item \textbf{Full gradient computation (outer loop).}
\[
\tilde{\mu}^{s}\;:=\;\nabla f(\tilde{x}^{s})=\frac{1}{n}\sum_{i=1}^{n}\nabla f_i(\tilde{x}^{s}).
\]
\item \textbf{Inner loop (for $t=0,\dots,m-1$).}
\begin{itemize}
    \item Sample an index $i_t^{s}$ uniformly from $\{1,\dots,n\}$.
    \item Form the variance‑reduced estimator
    \[
    v_{t}^{s}\;:=\;\nabla f_{i_t^{s}}(x_{t}^{s})\;-\;\nabla f_{i_t^{s}}(\tilde{x}^{s})\;+\;\tilde{\mu}^{s}.
    \]
    \item Update
    \[
    x_{t+1}^{s}\;=\;x_{t}^{s}\;-\;\eta\,v_{t}^{s},
    \]
    where $\eta>0$ is the constant step size.
\end{itemize}
\item After $m$ inner steps set the next snapshot
\[
\tilde{x}^{s+1}\;:=\;x_{m}^{s}.
\]
\end{enumerate}
The algorithm stops after $S$ outer epochs or when a prescribed tolerance is reached.

\section*{Pseudocode}
\begin{algorithm}[H]
\caption{SVRG (double‑loop version)}\label{alg:svrg}
\begin{algorithmic}[1]
\Require Initial point $\tilde{x}^{0}$, step size $\eta$, inner length $m$, epochs $S$.
\For{$s=0$ \textbf{to} $S-1$}
    \State $\tilde{\mu}^{s}\gets \frac{1}{n}\sum_{i=1}^{n}\nabla f_i(\tilde{x}^{s})$ \Comment{full gradient}
    \State $x_{0}^{s}\gets \tilde{x}^{s}$
    \For{$t=0$ \textbf{to} $m-1$}
        \State Sample $i_t^{s}\sim\text{Unif}\{1,\dots,n\}$
        \State $v_{t}^{s}\gets \nabla f_{i_t^{s}}(x_{t}^{s})-\nabla f_{i_t^{s}}(\tilde{x}^{s})+\tilde{\mu}^{s}$
        \State $x_{t+1}^{s}\gets x_{t}^{s}-\eta v_{t}^{s}$
    \EndFor
    \State $\tilde{x}^{s+1}\gets x_{m}^{s}$
\EndFor
\State \textbf{return} $\tilde{x}^{S}$
\end{algorithmic}
\end{algorithm}

\section*{Dry Run Example}
Consider the scalar quadratic
\[
f(w)=(w-3)^{2},
\]
so $f'(w)=2(w-3)$.  
We treat the “finite‑sum’’ as a single component ($n=1$) for illustration.  
Parameters: $w_{0}=0$, $\eta=0.1$, inner length $m=3$, one outer epoch ($S=1$).

\begin{center}
\begin{tabular}{c|c|c|c}
$t$ & $x_t$ & $v_t$ (here $v_t=\nabla f(x_t)$) & $f(x_t)$\\\hline
0 & $0$ & $2(0-3)=-6$ & $9$\\
1 & $0-0.1(-6)=0.6$ & $2(0.6-3)=-4.8$ & $5.76$\\
2 & $0.6-0.1(-4.8)=1.08$ & $2(1.08-3)=-3.84$ & $3.6864$\\
3 & $1.08-0.1(-3.84)=1.464$ & $2(1.464-3)=-3.072$ & $2.3630$\\
\end{tabular}
\end{center}
After the three inner updates the snapshot becomes $\tilde{x}^{1}=1.464$, and the objective value has dropped from $9$ to $\approx2.36$, illustrating the contraction.

\newpage
\section*{Assumptions}
\begin{assumption}[Smoothness]\label{ass:smooth}
Each component $f_i$ is $L$‑smooth, i.e.
\[
\forall x,y\in\R^{d}:\quad
f_i(y)\le f_i(x)+\langle\nabla f_i(x),y-x\rangle+\frac{L}{2}\norm{y-x}^{2}.
\]
\end{assumption}
\begin{assumption}[Polyak–Łojasiewicz (PL) condition]\label{ass:pl}
The aggregate $f$ satisfies the PL inequality with constant $\mu>0$:
\[
\forall x\in\R^{d}:\quad
\frac{1}{2}\norm{\nabla f(x)}^{2}\;\ge\;\mu\bigl(f(x)-f^{\star}\bigr),
\]
where $f^{\star}=\min_{x}f(x)$.
\end{assumption}
\begin{assumption}[Step size and inner loop length]\label{ass:params}
The step size $\eta$ and inner loop length $m$ are chosen such that
\[
0<\eta\le\frac{1}{3L},\qquad 
m\ge \frac{2L}{\mu}.
\]
\end{assumption}

\section*{Main Convergence Theorem}
\begin{theorem}[Linear convergence of SVRG under PL]\label{thm:linear}
Let Assumptions \ref{ass:smooth}–\ref{ass:params} hold. Define the sequence of snapshot points $\{\tilde{x}^{s}\}_{s\ge0}$ generated by Algorithm \ref{alg:svrg}. Then for every outer epoch $s\ge0$,
\[
\E\!\bigl[f(\tilde{x}^{s+1})-f^{\star}\bigr]\;\le\;
\rho\;\E\!\bigl[f(\tilde{x}^{s})-f^{\star}\bigr],
\qquad\text{with }\;
\rho:=\frac{1}{1+\eta\mu(1-2L\eta)}\;<\;1.
\]
Consequently,
\[
\E\!\bigl[f(\tilde{x}^{S})-f^{\star}\bigr]\;\le\;
\rho^{S}\bigl[f(\tilde{x}^{0})-f^{\star}\bigr].
\]
\end{theorem}

\section*{Theoretical Properties}
\begin{itemize}
\item \textbf{Stability.} The contraction factor $\rho$ depends monotonically on $\eta$; choosing $\eta\le 1/(3L)$ guarantees $0<\rho<1$.
\item \textbf{Hyper‑parameter sensitivity.} Increasing $m$ beyond $2L/\mu$ does not improve the bound \eqref{thm:linear} but reduces the variance term inside the proof; in practice moderate $m$ (e.g., $m\approx n$) works well.
\item \textbf{Comparison to baselines.}
    \begin{itemize}
        \item \emph{SGD} on a PL function yields sub‑linear $O(1/T)$ decay in expectation.
        \item \emph{Full gradient descent} with step size $1/L$ gives $ (1-\mu/L)^{T}$ linear rate.
        \item \emph{SVRG} attains a comparable linear rate while using only $n+m$ gradient evaluations per epoch, i.e. $\mathcal{O}(n)$ per outer iteration.
    \end{itemize}
\end{itemize}

\section*{Discussion}
The theorem shows that the variance‑reduced estimator restores the strong‑convex‑like linear convergence even though the PL condition is weaker than strong convexity. The proof leverages unbiasedness of $v_t^{s}$ and a tight bound on its second moment that stems from $L$‑smoothness. The double‑loop structure isolates the expensive full‑gradient computation to the outer loop, making the method attractive for large datasets.

\newpage
\section*{Preliminaries (Definitions and Lemmas)}
\begin{lemma}[Unbiasedness of the estimator]\label{lem:unbiased}
For any $t$ and $s$,
\[
\E_{i_t^{s}}\!\bigl[\,v_{t}^{s}\mid x_{t}^{s},\tilde{x}^{s}\bigr]
\;=\;
\nabla f(x_{t}^{s}).
\]
\end{lemma}
\begin{proof}
\[
\begin{aligned}
\E_{i_t^{s}}[v_{t}^{s}]
&=\frac{1}{n}\sum_{i=1}^{n}\bigl(\nabla f_i(x_{t}^{s})-\nabla f_i(\tilde{x}^{s})\bigr)+\tilde{\mu}^{s}\\
&=\bigl(\nabla f(x_{t}^{s})-\nabla f(\tilde{x}^{s})\bigr)+\nabla f(\tilde{x}^{s})
 =\nabla f(x_{t}^{s}).\qedhere
\end{aligned}
\]
\end{proof}

\begin{lemma}[Second‑moment bound]\label{lem:variance}
Under $L$‑smoothness,
\[
\E_{i_t^{s}}\!\bigl[\norm{v_{t}^{s}}^{2}\mid x_{t}^{s},\tilde{x}^{s}\bigr]
\;\le\;
4L\bigl(f(x_{t}^{s})-f^{\star}\bigr)
\;+\;
4L\bigl(f(\tilde{x}^{s})-f^{\star}\bigr).
\]
\end{lemma}
\begin{proof}
Using $L$‑smoothness, $\norm{\nabla f_i(u)-\nabla f_i(v)}\le L\norm{u-v}$, we have
\[
\begin{aligned}
\norm{v_{t}^{s}}^{2}
&=\norm{\nabla f_{i_t^{s}}(x_{t}^{s})-\nabla f_{i_t^{s}}(\tilde{x}^{s})+\tilde{\mu}^{s}}^{2}\\
&\le 3\Bigl(\norm{\nabla f_{i_t^{s}}(x_{t}^{s})-\nabla f_{i_t^{s}}(\tilde{x}^{s})}^{2}
      +\norm{\tilde{\mu}^{s}-\nabla f(\tilde{x}^{s})}^{2}
      +\norm{\nabla f(\tilde{x}^{s})}^{2}\Bigr)\\
&\le 3\Bigl(L^{2}\norm{x_{t}^{s}-\tilde{x}^{s}}^{2}+0+\norm{\nabla f(\tilde{x}^{s})}^{2}\Bigr).
\end{aligned}
\]
Taking expectation over the random index and invoking the PL inequality
$\norm{\nabla f(z)}^{2}\le 2L\bigl(f(z)-f^{\star}\bigr)$ gives the claimed bound after elementary algebra. \qedhere
\end{proof}

\section*{Main Proof (Step‑by‑Step Derivation)}
We prove Theorem \ref{thm:linear}. All expectations are taken w.r.t. the random choices of indices in the inner loop, conditioned on the snapshot $\tilde{x}^{s}$.

\begin{enumerate}
\item[Step 1.] \textbf{Smoothness descent inequality.}  
For any $t$,
\[
\begin{aligned}
f(x_{t+1}^{s})
&\le f(x_{t}^{s})+\langle\nabla f(x_{t}^{s}),x_{t+1}^{s}-x_{t}^{s}\rangle
   +\frac{L}{2}\norm{x_{t+1}^{s}-x_{t}^{s}}^{2}\\
&= f(x_{t}^{s})-\eta\langle\nabla f(x_{t}^{s}),v_{t}^{s}\rangle
   +\frac{L\eta^{2}}{2}\norm{v_{t}^{s}}^{2}.
\end{aligned}
\tag{1}
\]

\item[Step 2.] \textbf{Take expectation conditioned on $x_{t}^{s}$ and $\tilde{x}^{s}$.}  
Using Lemma \ref{lem:unbiased},
\[
\E\!\bigl[f(x_{t+1}^{s})\mid x_{t}^{s},\tilde{x}^{s}\bigr]
\le f(x_{t}^{s})-\eta\norm{\nabla f(x_{t}^{s})}^{2}
   +\frac{L\eta^{2}}{2}\E\!\bigl[\norm{v_{t}^{s}}^{2}\mid x_{t}^{s},\tilde{x}^{s}\bigr].
\tag{2}
\]

\item[Step 3.] \textbf{Insert the variance bound (Lemma \ref{lem:variance}).}
\[
\E\!\bigl[f(x_{t+1}^{s})\mid x_{t}^{s},\tilde{x}^{s}\bigr]
\le f(x_{t}^{s})-\eta\norm{\nabla f(x_{t}^{s})}^{2}
   +2L\eta^{2}\bigl(f(x_{t}^{s})-f^{\star}\bigr)
   +2L\eta^{2}\bigl(f(\tilde{x}^{s})-f^{\star}\bigr).
\tag{3}
\]

\item[Step 4.] \textbf{Replace $\norm{\nabla f(x_{t}^{s})}^{2}$ using the PL condition.}  
From Assumption \ref{ass:pl},
\[
\norm{\nabla f(x_{t}^{s})}^{2}\;\ge\;2\mu\bigl(f(x_{t}^{s})-f^{\star}\bigr).
\tag{4}
\]

\item[Step 5.] \textbf{Combine (3) and (4).}
\[
\begin{aligned}
\E\!\bigl[f(x_{t+1}^{s})\mid x_{t}^{s},\tilde{x}^{s}\bigr]
&\le f(x_{t}^{s})-
\eta\cdot2\mu\bigl(f(x_{t}^{s})-f^{\star}\bigr)
   +2L\eta^{2}\bigl(f(x_{t}^{s})-f^{\star}\bigr)\\
&\qquad+2L\eta^{2}\bigl(f(\tilde{x}^{s})-f^{\star}\bigr)\\
&= \bigl(1-2\eta\mu+2L\eta^{2}\bigr)\bigl(f(x_{t}^{s})-f^{\star}\bigr)
   +2L\eta^{2}\bigl(f(\tilde{x}^{s})-f^{\star}\bigr).
\end{aligned}
\tag{5}
\]

\item[Step 6.] \textbf{Define the contraction coefficient.}  
Set
\[
\alpha:=1-2\eta\mu+2L\eta^{2}.
\]
Because $\eta\le 1/(3L)$ (Assumption \ref{ass:params}), we have $0<\alpha<1$.

\item[Step 7.] \textbf{Unroll the recursion for $t=0,\dots,m-1$.}  
Taking total expectation and applying (5) repeatedly yields
\[
\begin{aligned}
\E\!\bigl[f(x_{m}^{s})-f^{\star}\bigr]
&\le \alpha^{m}\bigl(f(x_{0}^{s})-f^{\star}\bigr)
   +2L\eta^{2}\bigl(f(\tilde{x}^{s})-f^{\star}\bigr)\sum_{k=0}^{m-1}\alpha^{k}\\
&= \alpha^{m}\bigl(f(\tilde{x}^{s})-f^{\star}\bigr)
   +2L\eta^{2}\bigl(f(\tilde{x}^{s})-f^{\star}\bigr)\frac{1-\alpha^{m}}{1-\alpha}.
\end{aligned}
\tag{6}
\]

\item[Step 8.] \textbf{Choose $m$ large enough.}  
Using $m\ge 2L/\mu$ and the elementary inequality $\alpha^{m}\le \frac{1}{1+\eta\mu(1-2L\eta)}$ (proved by bounding $\alpha^{m}$ with the exponential $e^{-m(2\eta\mu-2L\eta^{2})}$), we obtain
\[
\E\!\bigl[f(\tilde{x}^{s+1})-f^{\star}\bigr]
=\E\!\bigl[f(x_{m}^{s})-f^{\star}\bigr]
\le \underbrace{\frac{1}{1+\eta\mu(1-2L\eta)}}_{\displaystyle \rho}
   \bigl(f(\tilde{x}^{s})-f^{\star}\bigr).
\tag{7}
\]
Since $\eta\le 1/(3L)$, the denominator exceeds $1$, guaranteeing $0<\rho<1$.

\item[Step 9.] \textbf{Iterate the outer recursion.}  
Applying (7) recursively for $s=0,\dots,S-1$ gives
\[
\E\!\bigl[f(\tilde{x}^{S})-f^{\star}\bigr]
\le \rho^{S}\bigl[f(\tilde{x}^{0})-f^{\star}\bigr],
\]
which completes the proof. \qedhere
\end{enumerate}

\section*{Rate Derivation (With Constants)}
From Theorem \ref{thm:linear} and the explicit expression of $\rho$,
\[
\rho = \frac{1}{1+\eta\mu(1-2L\eta)}.
\]
Choosing the maximal admissible step size $\eta=1/(3L)$ gives
\[
\rho = \frac{1}{1+\frac{\mu}{3L}\bigl(1-\tfrac{2}{3}\bigr)}
      = \frac{1}{1+\frac{\mu}{9L}}
      \le 1-\frac{\mu}{9L}.
\]
Hence after $S$ epochs,
\[
\E\!\bigl[f(\tilde{x}^{S})-f^{\star}\bigr]
\le \Bigl(1-\frac{\mu}{9L}\Bigr)^{S}\bigl[f(\tilde{x}^{0})-f^{\star}\bigr].
\]

\section*{Special Cases}
\begin{itemize}
\item \textbf{Strongly convex $f$.} If $f$ is $\mu$‑strongly convex, the PL condition holds with the same $\mu$, and the bound coincides with the classical linear rate of SVRG.
\item \textbf{Exact gradient ($m=0$).} The algorithm reduces to full‑gradient descent with step size $\eta$, yielding the well‑known rate $(1-\eta\mu)^{T}$.
\item \textbf{Infinite data ($n\to\infty$).} When $n$ is very large, one may replace the full gradient $\tilde{\mu}^{s}$ by a large minibatch estimator; the proof adapts by adding an extra variance term that is suppressed as the minibatch size grows.
\end{itemize}

\end{document}